\documentclass[a4paper,11pt]{article}
\usepackage[utf8]{inputenc}
\usepackage[T1]{fontenc}
\usepackage[french]{babel}
\usepackage{lmodern}
\usepackage{amsmath,amssymb,amsthm}
\usepackage{mathrsfs}
\usepackage{enumitem}
\usepackage{multicol}

\setlist[enumerate]{label=\textbf{\textcolor{Couleur8}{\arabic*.}}, left=1em}

% Si vous voulez aussi modifier les listes enumerate imbriquées :
\setlist[enumerate,2]{label=\textcolor{Couleur8}{\alph*)}}
\setlist[enumerate,3]{label=\textcolor{Couleur8}{\roman*)}}
\setlist[enumerate,4]{label=\textcolor{Couleur8}{\Alph*)}}
\usepackage{graphicx}
\usepackage{xcolor}
\usepackage{tcolorbox}
\usepackage{geometry}
\usepackage{fancyhdr}
\usepackage{titlesec}
\usepackage{cancel}
\usepackage{ifthen}
\usepackage{tikz}
\usetikzlibrary{arrows.meta}
\usepackage{tkz-tab}
\usepackage{eurosym}

% OPTION PROF/ÉLÈVE
% Décommentez UNE des deux lignes suivantes :
\newboolean{prof}\setboolean{prof}{true}   % Version professeur (avec corrections des devoirs)
\newboolean{prof}\setboolean{prof}{true}  % Version élève (sans corrections des devoirs)

\definecolor{Couleur1}{HTML}{4A5D7A} %Th
\definecolor{Couleur2}{HTML}{A5B5D6} %Prop
\definecolor{Couleur3}{HTML}{A8C68A} %Méthode
\definecolor{Couleur6}{HTML}{8A9CC1} %Def
\definecolor{Couleur7}{HTML}{E6B459} %Rq Voc
\definecolor{Couleur8}{HTML}{D36740} %Titres
\definecolor{correction}{HTML}{D36740} % Couleur pour les corrections
\definecolor{coopmathscorr}{HTML}{F15929} % Couleur corrections coopmaths

% Configuration des marges
\geometry{
	top=2cm,
	bottom=2cm,
	inner=1.5cm,
	outer=1.5cm,
}

% Police
\renewcommand{\familydefault}{\sfdefault}

% Compteurs
\newcounter{seancenum}
\newcounter{seancenumD}
\setcounter{seancenum}{1}
\setcounter{seancenumD}{1}

% Configuration des en-têtes et pieds de page
\pagestyle{fancy}
\fancyhf{}
\ifthenelse{\boolean{prof}}{
    \fancyhead[L]{\textcolor{Couleur8}{\textbf{Corrections Automatismes - Classe de seconde - Version Professeur}}}
}{
    \fancyhead[L]{\textcolor{Couleur8}{\textbf{Corrections Devoirs - Séance 1 - Classe de seconde}}}
}
\fancyhead[R]{\textcolor{Couleur8}{\textbf{2025-2026}}}
\fancyfoot[L]{\textcolor{Couleur8}{Mme Bertail et Mme Pinto}}
\fancyfoot[R]{\textcolor{Couleur8}{\thepage}}
\renewcommand{\headrulewidth}{1pt}
\renewcommand{\headrule}{\hbox to\headwidth{\color{Couleur8}\leaders\hrule height \headrulewidth\hfill}}

% Environnements personnalisés
\newtcolorbox{seance}[1]{
    colback=Couleur6!10,
    colframe=Couleur1!65,
    fonttitle=\bfseries\large,
    title=Correction Séance \theseancenum #1,
    left=5mm,
    right=5mm,
    top=3mm,
    bottom=3mm,
    arc=0mm,
    after=\stepcounter{seancenum}
}

\newtcolorbox{seanceD}[1]{
    colback=Couleur6!5,
    colframe=Couleur6!45,
    fonttitle=\bfseries\large,
    title=Correction Devoirs - Séance \theseancenumD #1,
    left=5mm,
    right=5mm,
    top=3mm,
    bottom=3mm,
    arc=0mm,
    after=\stepcounter{seancenumD}
}

\begin{document}

\setcounter{seancenumD}{1}
\newpage
\begin{seanceD}{} %1

\begin{enumerate}
\item $\dfrac{9}{2}+\dfrac{4}{20}=\dfrac{9{\color{Couleur1}\boldsymbol{\times 10}}}{2{\color{Couleur1}\boldsymbol{\times 10}}}+\dfrac{4}{20}=\dfrac{90+4}{20}=\dfrac{94}{20}=\dfrac{47\times2}{10\times2}=$ ${\color{correction}\boldsymbol{\dfrac{47}{10}}}$

On réduit au même dénominateur (20), puis on simplifie par 2.

\item $\dfrac{2}{6}-\dfrac{2}{9}=\dfrac{2{\color{Couleur1}\boldsymbol{\times 3}}}{6{\color{Couleur1}\boldsymbol{\times 3}}}-\dfrac{2{\color{Couleur1}\boldsymbol{\times 2}}}{9{\color{Couleur1}\boldsymbol{\times 2}}}=\dfrac{6-4}{18}=\dfrac{2}{18}=\dfrac{1\times2}{9\times2}=$ ${\color{correction}\boldsymbol{\dfrac{1}{9}}}$

Le dénominateur commun est 18, puis on simplifie par 2.

\item $\dfrac{\dfrac{14}{9}}{\dfrac{-63}{12}} = \dfrac{14}{9} \times \dfrac{12}{-63} = \dfrac{14 \times 12}{9 \times (-63)}$

On décompose : $\dfrac{2 \times 7 \times 3 \times 4}{9 \times (-9) \times 7} = \dfrac{2 \times \cancel{7} \times 3 \times 4}{9 \times (-9) \times \cancel{7}} = \dfrac{24}{-81} = {\color{correction}\boldsymbol{-\dfrac{8}{27}}}$

\item $\dfrac{14}{9}\times\dfrac{-3}{-35}=\dfrac{14\times(-3)}{9\times(-35)} = \dfrac{-42}{-315} = \dfrac{42}{315}$

On décompose : $\dfrac{2 \times 3 \times 7}{9 \times 5 \times 7} = \dfrac{2 \times 3 \times \cancel{7}}{9 \times 5 \times \cancel{7}} = \dfrac{6}{45} = {\color{correction}\boldsymbol{\dfrac{2}{15}}}$

\item Le matin, Fiona a bu $\dfrac{1}{3}$ de la bouteille. Il reste alors $1 - \dfrac{1}{3} = \dfrac{2}{3}$ de la bouteille.

À midi, elle a bu $\dfrac{1}{4}$ du reste, soit : $\dfrac{1}{4} \times \dfrac{2}{3} = {\color{correction}\boldsymbol{\dfrac{2}{12} = \dfrac{1}{6}}}$ de la bouteille.

\item $A= \dfrac{3}{4-\dfrac{6}{7}} = \dfrac{3}{\dfrac{28}{7}-\dfrac{6}{7}} = \dfrac{3}{\dfrac{22}{7}} = 3 \times \dfrac{7}{22} = {\color{correction}\boldsymbol{\dfrac{21}{22}}}$

\end{enumerate}

\end{seanceD}

\end{document}